\documentclass[a4paper, 11pt,twoside=false]{scrartcl} % Remplace la classe "article", correspond au standard européen.
\usepackage{Rapport2}
\usepackage{ctex}
\usepackage{indentfirst}
\setCJKmainfont{思源宋体}
%\usepackage[toc]{multitoc} double table des matières si nécessaire
\usepackage{lipsum}
\clearpairofpagestyles


%%% Paramètres pour les footers et la page de titre
\newcommand{\lab}{计算书}
\newcommand{\tit}{电动力绳弹射装置设计}
\newcommand{\nom}{}
\newcommand{\superv}{}
\newcommand{\depa}{}

\automark[section]{section}

\lehead*{%
  \makebox[0pt][r]{
    \colorbox{rahmen}{\makebox[\textwidth][r]{\bfseries\color{white}{\pagemark/\pageref*{LastPage}}\enskip}}\hspace*{1em}}%
  \headmark\hspace*{1em}\headrulefill%
}
\rohead*{%
  \headrulefill\hspace*{1em}\headmark%
  \makebox[0pt][l]{%
    \hspace*{1em}\colorbox{rahmen}{\makebox[\textwidth][l]{\enskip\bfseries\color{white}{\pagemark/\pageref*{LastPage}}}}}%
}
\lefoot{\color{darkgray}\small \today}
\rofoot{\color{darkgray}\small \today}
\cefoot{\color{darkgray}\small \itshape \tit\unskip\strut}
\cofoot{\color{darkgray}\small \itshape \tit\unskip\strut}
\refoot{\color{darkgray}\small \nom\unskip\strut}
\lofoot{\color{darkgray}\small \nom\unskip\strut}

\addtokomafont{pagenumber}{\bfseries\color{white}}
\newcommand\headrulefill{\leaders\hrule width 0pt height 3pt depth -2.8pt \hfill}

%\renewcommand\titlepagestyle{empty}

\hfuzz 100pt
\hbadness 10000
\hfuzz 100pt
\hbadness 10000

% Inversion des espaces des marges

\let\tempmargin\oddsidemargin
\let\oddsidemargin\evensidemargin
\let\evensidemargin\tempmargin
\reversemarginpar

\sisetup{detect-weight=true, detect-family=true}

\renewcommand*{\maketitle}{%
\begin{titlepage}
\raggedleft\includegraphics[scale=0.5]{sect_physique_pant}

\begin{center}
\horrule{0.5pt} \\[0.4cm] \vspace{-1.5ex}\textcolor{darkblue}
	 { \huge \bfseries \lab \\ \vspace{0.4cm} \tit }\\[0.1cm]\horrule{2pt} \vspace{-2ex}
\end{center}
\vspace{1cm}
\begin{center}
 \textsc{\LARGE 北京理工大学\\ \large 宇航学院 \\ \depa\unskip\strut} \vspace{-1.5ex} \\	
\end{center}
\vspace{3cm}
\begin{minipage}[t]{0.5\textwidth}
    \begin{flushleft}
        {\large \nom\unskip\strut
        }
    \end{flushleft}
\end{minipage}%
%
\begin{minipage}[t]{0.5\textwidth}
    \begin{flushright}
        {\large \textit{导师}\superv\unskip\strut}           
        \vspace*{0.5cm}
    \end{flushright}    
\end{minipage}%

\vspace{4cm}
\begin{center}
\begin{minipage}[b]{0.5\textwidth}
    \vfill
    \begin{center}
        {\large\today}
    \end{center}    
\end{minipage}%
\end{center}

\end{titlepage}
}




\begin{document}


%%%%TITRE

\maketitle

\tableofcontents % table des matières
\newpage

%%%%%%% Début du document

\section{迁移}
她叫林徽因，出生于杭州，是许多人梦中期待的白莲。她在雨雾之都伦敦，
发生过一场空前绝后的康桥之恋。她爱过三个男子，爱得清醒，也爱得平静。徐
志摩为她徜徉在康桥，深情地等待一场旧梦可以归来。梁思成与她携手走过千山
万水，为完成使命而相约白头。金岳霖为她终身不娶，痴心不改地守候一世。可
她懂得人生飘忽不定，要学会随遇而安。真正的平静，不是避开车马喧嚣，而是
在心中修篱种菊。尽管如流往事，每一天都涛声依旧，只要我们消除执念，便可
寂静安然。愿每个人在纷呈世相中不会迷失荒径，可以端坐磐石上，醉倒落花
前。如果可以，请让我预支一段如莲的时光，哪怕将来某一天加倍偿还。这个雨
季会在何时停歇，无从知晓。但我知道，你若安好，便是晴天。
\section{公式}
Soit $f:E\to F$ une fonction  dans $E$. Pour $x\in E$, on a
\[
\int_0^x f'(y)\dd{y}=f(x)-f(0).
\]
Arrays of mathematics are typeset using one of the matrix environments as 
in
\[
\begin{bmatrix}
1 & x & 0 \\
0 & 1 & -1
\end{bmatrix}\begin{bmatrix}
1  \\
y  \\
1
\end{bmatrix}
=\begin{bmatrix}
1+xy  \\
y-1
\end{bmatrix}.
\]
Case statements use cases:
\[
|x|=\begin{cases}
x, & \text{if }x\geq 0\,,  \\
-x, & \text{if }x< 0\,.
\end{cases}
\]
Many arrays have lots of dots all over the place as in
\[
\begin{matrix}
-2 & 1 & 0 & 0 & \cdots & 0  \\
1 & -2 & 1 & 0 & \cdots & 0  \\
0 & 1 & -2 & 1 & \cdots & 0  \\
0 & 0 & 1 & -2 & \ddots & \vdots \\
\vdots & \vdots & \vdots & \ddots & \ddots & 1  \\
0 & 0 & 0 & \cdots & 1 & -2
\end{matrix}
\]


\section{不错的话}
她叫林徽因，出生于杭州，是许多人梦中期待的白莲。她在雨雾之都伦敦，
发生过一场空前绝后的康桥之恋。她爱过三个男子，爱得清醒，也爱得平静。徐
志摩为她徜徉在康桥，深情地等待一场旧梦可以归来。梁思成与她携手走过千山
万水，为完成使命而相约白头。金岳霖为她终身不娶，痴心不改地守候一世。可
她懂得人生飘忽不定，要学会随遇而安。真正的平静，不是避开车马喧嚣，而是
在心中修篱种菊。尽管如流往事，每一天都涛声依旧，只要我们消除执念，便可
寂静安然。愿每个人在纷呈世相中不会迷失荒径，可以端坐磐石上，醉倒落花
前。如果可以，请让我预支一段如莲的时光，哪怕将来某一天加倍偿还。这个雨
季会在何时停歇，无从知晓。但我知道，你若安好，便是晴天。
\section{ et analyse}
\lipsum[30]
\section{Conclusion}
\lipsum[20]
\end{document}
